\section*{Phase 1: Data Science Project Proposals}

\subsection*{Project 1: Housing Amenities Impact on Apartment Listing Prices in the Southern USA}

\subsubsection*{Business Issue and Overview}

Real estate investors and property managers in the Southern USA lack quantitative evidence on which amenities drive rental pricing. Understanding amenity importance enables strategic property improvements and competitive pricing. This project analyzes which amenities (parking, laundry, pet policies, fitness centers, etc.) most significantly influence rental prices in major Southern markets: Austin TX, Dallas TX, Miami FL, Atlanta GA, Nashville TN, and Charlotte NC. Geographic focus is essential: amenity value in Austin differs significantly from Miami or Atlanta, requiring region-specific analysis.

\subsubsection*{Business Questions}

\begin{itemize} [itemsep = -9pt]
    \item Which amenities have strongest impact on rental pricing?
    \item How does amenity importance vary across geographic regions?
    \item Which amenity upgrades yield the highest price premium relative to cost?
    \item Do certain amenity bundles (e.g., parking \& laundry) have compounding effects on price?
\end{itemize}

\subsubsection*{Data Needs}

Major entities: Rental listings from Southern markets (Austin, Dallas, Miami, Atlanta, Nashville, Charlotte) including price, location, bedrooms, bathrooms, square footage, amenities (binary/categorical flags), listing metadata (date listed, availability). Requires at least two datasets for robust analysis (e.g., Zillow and Apartments.com data combined, or regional MLS databases).

\subsubsection*{Reason}

Addresses urgent business problem in high-growth Southern real estate markets. Understanding which amenities justify investment is directly actionable for property managers operating in Texas, Florida, and Carolina markets. Regional focus ensures recommendations are relevant to Southern investors rather than generic national insights. Deepens familiarity with feature importance techniques and regression modeling while producing region-specific competitive insights.

\newpage

\subsection*{Project 2: Student Course Success Prediction in Southern Universities}

\subsubsection*{Business Issue and Overview}

Educational institutions across the Southern USA (focus: large state universities in Texas, Georgia, Florida) face resource constraints and must identify at-risk students early. This project builds a predictive model to identify undergraduate students likely to struggle or fail in introductory STEM courses, enabling early intervention programs. Predictions target first-semester students enrolled at these Southern institutions to maximize intervention effectiveness and improve retention rates in competitive regional markets.

Which student characteristics predict success in STEM courses at Southern universities?
\begin{itemize} [itemsep = -9pt]
    \item How do student characteristics vary by institution type (large state vs. smaller regional)?
    \item What baseline high school GPA and SAT/ACT scores indicate risk for Southern STEM courses?
    \item When in the semester should interventions occur for maximum effectiveness?
    \item How does geographic region (in-state vs. out-of-state) impact success prediction?
    \item How accurately can we identify at-risk students for early warning systems?
\end{itemize}

\subsubsection*{Data Needs}

Major entities: Student demographics (age, prior GPA, SAT/ACT scores, high school region, in-state/out-of-state status, major), school characteristics (university name, location in Southern region, institution type, enrollment size), course characteristics (STEM subject, difficulty level, instructor experience, class size), performance metrics (grades, attendance, engagement, pass/fail status). Minimum two datasets required: one from university registrar systems and one from institutional research databases or secondary education datasets.

\subsubsection*{Reason}

Addresses critical institutional challenge in Southern universities: student retention and success in competitive STEM programs. Explicit focus on undergraduate students at Southern institutions (Texas, Georgia, Florida) provides actionable insights for regional education stakeholders. Classification modeling provides interpretable insights; understanding why students succeed in STEM courses enables targeted, evidence-based intervention strategies. Valuable for university administration, advising departments, and academic support centers across the South. Demonstrates practical impact of data science on human outcomes and institutional effectiveness while generating region-specific recommendations.

%\subsection*{Recommended Project}

%\textbf{Project 1 (Housing Amenities)} is recommended. It aligns with your prior work, has established data availability (Craigslist listings), and delivers concrete business value for real estate stakeholders. The feature importance analysis directly answers investor questions about amenity ROI.
